\documentclass[10pt]{beamer}

\usetheme{metropolis}
\definecolor{WiLabRed}{RGB}{197,18,48}
\setbeamercolor{frametitle}{fg=white,bg=WiLabRed}
\setbeamercolor{progress bar}{fg=WiLabRed!90}
\setbeamercolor{title separator}{fg=WiLabRed!90}
\setbeamercolor{progress bar in section page}{fg=WiLabRed!90}
\setbeamercolor{background canvas}{bg=white}
\setbeamercolor{alerted text}{fg=WiLabRed!90}

\usepackage{appendixnumberbeamer}

\usepackage{booktabs}
\usepackage[scale=2]{ccicons}

\usepackage{pgfplots}
\usepgfplotslibrary{dateplot}

\usepackage{xspace}
\newcommand{\themename}{\textbf{\textsc{metropolis}}\xspace}

\usepackage{marvosym}
%\usepackage{subfig}
\usepackage{graphicx}\graphicspath{{images/}}
\usepackage{subcaption}
\usepackage[framed]{mcode}
\usepackage{listings}
\usepackage{amsmath}

\DeclareMathOperator{\sinc}{sinc}
\DeclareMathOperator{\sgn}{sgn}

\title{Lecture 13: Fine Frequency Correction}
\subtitle{\textit{Software Defined Radio for Engineers} (Collins~\textit{et al.}), \textsection{7.2.2}}
\date{}
\author{\textbf{Alexander M. Wyglinski, Ph.D.}}
\institute{ \vspace*{1in}\hfill\includegraphics[height=1.125cm]{wilab_logo-A70916.eps} \qquad \includegraphics[height=1.125cm]{WPI_Inst_Prim_FulClr.eps}}
% \titlegraphic{\hfill\includegraphics[height=1.5cm]{logo.pdf}}

% Foot for all slides
\setbeamertemplate{frame footer}{\tiny \copyright~2018 by Alexander Wyglinski. This work is licensed under the Creative Commons Attribution-ShareAlike 4.0 International License. To view a copy of this license, visit http://creativecommons.org/licenses/by-sa/4.0/.}

\begin{document}

%\captionsetup[subfigure]{labelformat=empty}

%%%%%%%%%%%%%%%%%%%%%%%%%%%%%%%%%%%%%%%%%%%%%%%%%%%%%%%%%%

\maketitle




%%%%%%%%%%%%%%%%%%%%%%%%%%%%%%%%%%%%%%%%%%%%%%%%%%%%%%%%%%

\frame
{
  \frametitle{FFC Purpose}

    \begin{itemize}
        \item Once Coarse Frequency Correction (CFC) is performed, Fine Frequency Correction (FFC) is applied to identify exact location of carrier frequency
	\begin{itemize}
	 \item Some digital signals might be susceptible to small amounts of carrier frequency offset $\rightarrow$ Might cause phase rotation in signal constellation
	 \item Rate of rotate equal to frequency offset
	\end{itemize}
        \item Unlike CFC, which uses feedforward approach, FFC uses feedback approach to find carrier frequency
        \begin{itemize}
         \item Need to implement digital phase locked loop (DPLL)
        \end{itemize}
    \end{itemize}
}


%%%%%%%%%%%%%%%%%%%%%%%%%%%%%%%%%%%%%%%%%%%%%%%%%%%%%%%%%%

\frame
{
  \frametitle{FFC DPLL Overview}

\centering
$x(t)$\\
$\downarrow$\\
$\phi(n)$ $\rightarrow$ Phase Rotator $\rightarrow$ $y(n)$\\
$\downarrow$\\
Phase Error Detector $\rightarrow$ $e(n)$\\
$\downarrow$\\
Loop Filter $\rightarrow$ $f(n)$\\
$\downarrow$\\
Direct Digital Synthesizer $\rightarrow$ $\phi(n)$\\
  
}


%%%%%%%%%%%%%%%%%%%%%%%%%%%%%%%%%%%%%%%%%%%%%%%%%%%%%%%%%%

\frame
{
  \frametitle{FFC DPLL Variables}
 
 \begin{itemize}
  \item $x(t)$: Input signal to DPLL that possesses some degree of carrier offset
  \item $y(n)$: Corrected signal with respect to carrier offset
  \item $e(n)$: Error signal produced by phase error detector (PED), defined by:
  \begin{equation}
   e(n)=\sgn(\Re(y(n)))\cdot\Im(y(n))-\sgn(\Im(y(n)))\cdot\Re(y(n))
  \end{equation}
  which is essentially the difference between the real and imaginary components of $y(n)$ and will only be zero when $\Re(y(n))=\Im(y(n))$
  \item $f(n)$: Output of loop filter, which governs dynamics of DPLL by determining the operational frequency, lock time, and responsiveness (also smooths out error)
  \item $\phi(n)$: The correction signal that is applied to input signal $x(t)$ to yield compensated output signal $y(n)$
 \end{itemize}

 
}

%%%%%%%%%%%%%%%%%%%%%%%%%%%%%%%%%%%%%%%%%%%%%%%%%%%%%%%%%%

\frame
{
  \frametitle{Loop Filter Design}

  \begin{itemize}
   \item Once we have $e[n]$, the next step is to apply the loop filter
   \item Calculations of gain values $(G_1,G_2)$ employ the following expressions:
\begin{align}
 \theta=\frac{B_{\mathrm{Loop}}}{M(\psi+0.25/\psi)},\quad\Delta=1+2\psi\theta+\psi^2\\
 G_1=\frac{4\psi\theta/\Delta}{MK},\quad{G_2}=\frac{(4/M)\theta^2/\Delta}{MK}
\end{align}
such that $B_{\mathrm{Loop}}$ is loop filter bandwidth, $\psi$ is preferred dampening factor, $M$ is number of samples per second, and $K$ is detector gain
\item: Helpful Hint: When designing a loop filter, start off with $B_{\mathrm{Loop}}=0.01$ and $\psi=1$ since it is easier to work with an overdampened system
  \end{itemize}
 
}

%%%%%%%%%%%%%%%%%%%%%%%%%%%%%%%%%%%%%%%%%%%%%%%%%%%%%%%%%%


\frame
{
  \frametitle{Direct Digital Synthesizer}

  \begin{itemize}
   \item The DDC is equivalent to an integrator, namely:
\begin{equation}
 F(z)=G_3\frac{z^{-1}}{1-z^{-1}}
\end{equation}
\end{itemize}

}

\end{document}
